\documentclass[11pt,a4paper]{article}

\usepackage{fontspec}
\setmainfont{DejaVu Serif}[
  BoldFont={DejaVu Serif Bold},
  ItalicFont={DejaVu Serif},
  ItalicFeatures={FakeSlant=0.18},
  BoldItalicFont={DejaVu Serif Bold},
  BoldItalicFeatures={FakeSlant=0.18}
]
\setsansfont{DejaVu Sans}
\usepackage{polyglossia}
\setmainlanguage{english}
\usepackage[a4paper,margin=20.5mm]{geometry}
\usepackage{microtype}
\usepackage{setspace}
\setstretch{1.035}
\usepackage{amsmath,amssymb,amsthm,mathtools,bm}
\usepackage{booktabs,longtable,array,tabularx}
\usepackage{enumitem}
\usepackage{xcolor}
\usepackage{tikz}
\usetikzlibrary{arrows.meta,positioning,calc}
\usepackage[unicode,hidelinks]{hyperref}
\usepackage[nameinlink,noabbrev]{cleveref}
\emergencystretch=3em
\allowdisplaybreaks

\hypersetup{
  pdftitle={Relativistic Beam Paths as Observation Geometry},
  pdfauthor={Stas, Independent Research Program},
  pdfsubject={Typed LHC beam kinematics, observation channels, accelerator invariants, and detector reconstruction},
  pdfkeywords={LHC, accelerator physics, relativistic kinematics, observation geometry, natural units, emittance, invariant mass}
}

\definecolor{obsblue}{RGB}{37,83,138}
\definecolor{obsred}{RGB}{150,45,45}
\definecolor{obsgreen}{RGB}{30,115,78}
\definecolor{lightblue}{RGB}{239,246,253}
\definecolor{lightred}{RGB}{253,242,242}
\definecolor{lightgreen}{RGB}{239,249,244}

\theoremstyle{definition}
\newtheorem{definition}{Definition}[section]
\newtheorem{model}[definition]{Model}
\theoremstyle{plain}
\newtheorem{proposition}[definition]{Proposition}
\theoremstyle{remark}
\newtheorem{remark}[definition]{Remark}
\newtheorem{warning}[definition]{Boundary}

\newcommand{\R}{\mathbb{R}}
\newcommand{\dd}{\mathop{}\!d}
\newcommand{\abs}[1]{\left\lvert #1\right\rvert}
\newcommand{\norm}[1]{\left\lVert #1\right\rVert}
\newcommand{\Law}{\operatorname{Law}}
\newcommand{\Cov}{\operatorname{Cov}}
\newcommand{\diag}{\operatorname{diag}}
\newcommand{\GeV}{\mathrm{GeV}}
\newcommand{\TeV}{\mathrm{TeV}}
\newcommand{\Ocal}{\mathcal O}
\newcommand{\Kcal}{\mathcal K}
\newcommand{\Rcal}{\mathcal R}
\newcommand{\Ecal}{\mathcal E}
\newcommand{\Fcal}{\mathcal F}
\newcolumntype{L}[1]{>{\raggedright\arraybackslash}p{#1}}
\newcolumntype{Y}{>{\raggedright\arraybackslash}X}

\title{\bfseries Relativistic Beam Paths as Observation Geometry\\
\large From synchrotron phase space to typed detector invariants}
\author{Stas\\
\small Independent Research Program; ORCID: 0009-0000-2294-705X}
\date{Strict GO Core reference revision v1.3 --- 28 July 2026}

\begin{document}
\maketitle

\begin{center}
\fcolorbox{obsblue}{lightblue}{%
\parbox{0.92\linewidth}{\small
\textbf{Status.} This document supersedes the bilingual working draft
v1.2. It is migrated to GO Core v0.2 and the SI--HEP Quantity Passport
v0.5. The physical reference is standard Run 3 proton--proton
operation; on the document date the LHC is in Long Shutdown 3. No new
accelerator model or particle-physics prediction is claimed.}}
\end{center}

\begin{abstract}
An accelerator is not observed as one visible trajectory. Its hidden
state contains a lattice-constrained reference orbit, time-dependent
beam distributions, RF and betatron phases, magnet settings, collision
states, and detector response variables. We type the corresponding
chain as
\[
X_t\xrightarrow{\Rcal}\widetilde X_t
\xrightarrow{\Ocal_\lambda}Z_t
\xrightarrow{\Kcal_\lambda}\Law(Y_t\mid Z_t)
\xrightarrow{\widehat\theta_\lambda}
\widehat\theta_t,
\]
where reduction, ideal observation, stochastic measurement, and
estimation are distinct maps. The SI--HEP boundary is explicit:
\(p^\mu_{\rm SI}=(E/c,\bm p)\),
\(\Pi^\mu=(E,c\bm p)\), and the HEP numerical four-vector is
\(\Pi^\mu/\Ecal_0\). This removes the ambiguity in
\(E^2=\bm p^2+m^2\). For charged beam transport, SI magnetic rigidity
and HEP momentum are connected by the exact declared-unit identity
\[
p[\GeV/c]=0.299792458\,z\,B\rho[\mathrm{T\,m}].
\]
At a Run 3 reference total energy of \(6.8\,\TeV\), the proton has
\(\gamma\approx7.2474\times10^3\),
\(1-\beta\approx9.5194\times10^{-9}\), and
\(B\rho\approx2.26824\times10^4\,\mathrm{T\,m}\). The manuscript
separates Lorentz invariants, longitudinal-boost covariants,
symplectic invariants, operational observables, and estimator outputs.
In particular, geometric emittance is invariant only under a specified
ideal symplectic transport, while normalized emittance is the relevant
ideal adiabatic invariant during acceleration. Equal-time Lorentz
contraction and Penrose--Terrell camera appearance are retained as
different observation maps, neither being a direct imaging model for
LHC protons.
\end{abstract}

\tableofcontents

\section{Scope and claim boundary}

The manuscript is a typed interface between standard accelerator
physics and Geometry of Observation. It is not:
\begin{itemize}
  \item a lattice design, collimation, impedance, beam-loss, or
  machine-protection model;
  \item a replacement for tracking, beam-beam, or detector simulation;
  \item a hard-body model of the proton;
  \item a detector reconstruction algorithm;
  \item a claim that a hadron's Lorentz-contracted equal-time profile
  is photographed in the LHC;
  \item a new Standard Model result.
\end{itemize}

The strict contribution is infrastructural: it identifies the objects,
maps, unit charts, transformation groups, and claim conditions needed
to state accelerator observations without type or dimension gaps.

\section{Observation protocol}

\begin{definition}[Hidden accelerator state]
A minimal hidden state at machine time \(t\) is a record
\[
X_t=
\bigl(
\Gamma_0,\mathcal L_{\rm mag}(t),\Phi_{\rm RF}(t),
f_1(t),f_2(t),\mathcal C_t,\mathcal D_t
\bigr),
\]
where \(\Gamma_0\) is a design orbit, \(\mathcal L_{\rm mag}\) is the
magnet lattice state, \(\Phi_{\rm RF}\) is the RF state,
\(f_1,f_2\) are beam phase-space distributions,
\(\mathcal C_t\) is a collision-event state when a crossing occurs,
and \(\mathcal D_t\) contains detector-response latent variables.
\end{definition}

An observation protocol \(\lambda\) contains a frame, clock,
calibration, spatial and temporal resolutions, trigger and selection
rules, a deterministic response, a stochastic channel, and an
estimator. Its typed chain is
\[
X_t\xrightarrow{\Rcal}\widetilde X_t
\xrightarrow{\Ocal_\lambda}Z_t
\xrightarrow{\Kcal_\lambda}Y_t
\xrightarrow{\widehat\theta_\lambda}\widehat\theta_t.
\]
Here \(\Kcal_\lambda\) denotes a Markov kernel. It must not be merged
with the deterministic map \(\Ocal_\lambda\), and
\(\widehat\theta_\lambda\) must not be identified with the underlying
measurand.

\begin{center}
\begin{tikzpicture}[
  node distance=8mm and 10mm,
  box/.style={draw,rounded corners,minimum width=29mm,minimum height=11mm,align=center},
  arr/.style={-{Latex[length=1.9mm]},thick}
]
\node[box,fill=lightblue] (x) {hidden state\\\(X_t\)};
\node[box,right=of x] (z) {ideal readout\\\(Z_t\)};
\node[box,right=of z] (y) {raw data\\\(Y_t\)};
\node[box,fill=lightgreen,below=of z] (est) {estimate\\\(\widehat\theta_t,\Sigma_t\)};
\draw[arr] (x) -- node[above]{\(\Rcal,\Ocal_\lambda\)} (z);
\draw[arr] (z) -- node[above]{\(\Kcal_\lambda\)} (y);
\draw[arr] (y) |- node[pos=0.72,right]{\(\widehat\theta_\lambda\)} (est);
\end{tikzpicture}
\end{center}

\section{Reference machine passport}

The numerical audit uses the CERN public Run 3 reference:
\[
E_{\rm beam}=6.8\,\TeV,\qquad
\sqrt{s}=13.6\,\TeV
\]
for equal, head-on proton beams. The public circumference
\(C=26659\,\mathrm m\), roughly \(2500\) bunches per proton beam,
and roughly \(1.6\times10^{11}\) protons per bunch are rounded or
approximate operational descriptors, not defining constants.

\begin{warning}[Epoch]
The LHC entered Long Shutdown 3 in June 2026. Therefore
\(6.8\,\TeV\) in this document is a Run 3 standard-operation
reference, not a statement that a beam was circulating on
28 July 2026.
\end{warning}

Every machine record must include
\[
(\mathsf{facility},\mathsf{species},\mathsf{charge\ state},
\mathsf{run},\mathsf{mode},\mathsf{epoch},\mathsf{status},
\mathsf{value},\mathsf{unit},\mathsf{precision},\mathsf{source}).
\]
The design energy, an achieved beam energy, a special low-energy run,
and an HL-LHC target are different records even when they share a unit.

\section{SI--HEP representation boundary}

Adopt \(\eta=\diag(1,-1,-1,-1)\). The SI momentum four-vector and
energy four-vector are
\[
p^\mu_{\rm SI}=(E/c,\bm p_{\rm SI}),\qquad
\Pi^\mu=c\,p^\mu_{\rm SI}=(E,c\bm p_{\rm SI}).
\]
With energy anchor \(\Ecal_0=1\,\GeV\), define the HEP coordinate
\[
P^\mu_{\rm HEP}=\frac{\Pi^\mu}{\Ecal_0}.
\]
The mass shell has three equivalent representations:
\[
p_{\rm SI}^2=m^2c^2,\qquad
\Pi^2=(mc^2)^2,\qquad
P_{\rm HEP}^2=\left(\frac{mc^2}{\Ecal_0}\right)^2.
\]

The shorthand \(E^2=\bm p^2+m^2\) is used only in the last chart.
In the other charts, omitted factors of \(c\) change dimensions.
The complete conversion contract is specified in the companion
\emph{SI--HEP Quantity Passport v0.5}.

\section{Relativistic beam kinematics}

Let
\[
\mu_p:=m_pc^2=0.93827208943(29)\,\GeV
\]
from CODATA 2022, and let \(E\) denote proton total energy. Define the
energy-valued momentum \(\pi:=pc\). Then
\[
E^2=\pi^2+\mu_p^2,\qquad
\gamma=\frac{E}{\mu_p},\qquad
\beta=\frac{\pi}{E}.
\]
At high \(\gamma\), direct subtraction of \(\beta\) from one loses
floating-point significance. The stable identity is
\[
1-\beta
=\frac{1-\beta^2}{1+\beta}
=\frac{\gamma^{-2}}{1+\beta}.
\]
This quantity, rather than a rounded decimal for \(\beta\), controls
the speed deficit:
\[
c-v=c(1-\beta).
\]

\begin{proposition}[Equal head-on beams]
For two equal beams with four-momenta
\(\Pi_1=(E,\bm\pi)\) and
\(\Pi_2=(E,-\bm\pi)\),
\[
s=(\Pi_1+\Pi_2)^2=4E^2,\qquad \sqrt{s}=2E.
\]
\end{proposition}

\begin{proof}
The total spatial momentum vanishes, while the total energy is \(2E\).
\end{proof}

The result assumes equal, exactly head-on beams and neglects event-wise
energy spread and crossing-angle corrections. Those effects belong to
the event and beam-distribution passports.

\section{Frame slices and light-cone observations}

For a boost of rapidity \(Y\) along the beam axis,
\[
\begin{aligned}
ct'&=\gamma(ct-\beta z),&
z'&=\gamma(z-\beta ct),\\
x'&=x,&y'&=y.
\end{aligned}
\]
The interval
\[
\Delta s^2=c^2\Delta t^2-\Delta x^2-\Delta y^2-\Delta z^2
\]
is Lorentz invariant. A rest-defined longitudinal proper length
\(L_0\), evaluated between endpoints simultaneous in the laboratory,
has laboratory equal-time length \(L_\parallel=L_0/\gamma\).
This is a relation between simultaneity slices, not a material
compression law for a quantum proton.

In light-front coordinates
\[
x^\pm=\frac{ct\pm z}{\sqrt2},
\]
the same boost acts as
\[
x^{+\prime}=e^{-Y}x^+,\qquad
x^{-\prime}=e^Yx^-.
\]

A camera image is a different map. A source worldline
\(X^\mu(\xi,\tau)\) contributes to an image at observer event
\(O^\mu\) only if
\[
\bigl(O-X(\xi,\tau_e)\bigr)^2=0,\qquad \tau_e<t_{\rm obs}.
\]
The incoming null direction is then projected to an image plane.
Thus:
\[
\text{equal-time boost}\ne\text{past-light-cone image}.
\]
The Penrose--Terrell effect concerns the second map. It is not a
detector-observable model for individual LHC protons.

\section{Magnetic rigidity and path geometry}

In SI,
\[
\frac{\dd\bm p}{\dd t}
=q(\bm E+\bm v\times\bm B),\qquad
\frac{\dd E}{\dd t}=q\,\bm E\cdot\bm v.
\]
A transverse magnetic field bends the orbit but does no work.
For local curvature radius \(\rho\),
\[
\abs{\bm p}_{\rm SI}=\abs q\,B\rho.
\]
With \(q=ze\), the exact declared-unit bridge is
\[
\boxed{
p[\GeV/c]=0.299792458\,z\,B\rho[\mathrm{T\,m}].
}
\]
For a Run 3 reference proton,
\[
p=6799.999935268\,\GeV/c,\qquad
B\rho=22682.358258\,\mathrm{T\,m}.
\]

\begin{warning}[Local versus global geometry]
The circumference contains straight insertions and does not determine
the dipole bending radius by \(C/(2\pi)\). The local field-to-curvature
map requires the lattice geometry. Conversely, the revolution
frequency may use the total closed-orbit length.
\end{warning}

\section{RF phase and longitudinal state}

For an idealized RF gap,
\[
\Delta E=qV_{\rm eff}\sin\phi_s.
\]
For \(q=ze\),
\[
\Delta E[\mathrm{eV}]
=zV_{\rm eff}[\mathrm V]\sin\phi_s.
\]
The unit conversion is exact; the effective-voltage and thin-gap
model is not.

A common dimensionless longitudinal coordinate is
\[
\delta=\frac{p-p_0}{p_0}.
\]
The phase \(\phi\) and \(\delta\) are not direct positions and momenta
until a canonical normalization is declared. Consequently, a plotted
RF ``bucket area'' is not comparable across conventions without the
Jacobian and reference energy. The observation protocol must record
the RF harmonic, synchronous phase, voltage convention, longitudinal
coordinate convention, and clock.

\section{Transverse optics and emittance}

In one transverse plane, linearized motion about the design orbit is
\[
x''(s)+K(s)x(s)=0,
\qquad [K]=L^{-2}.
\]
Here \(x\) has dimension \(L\), \(s\) has dimension \(L\), and
\(x'=\dd x/\dd s\) is dimensionless. A first-order transport map
\(\bm u_2=M\bm u_1\), \(\bm u=(x,x')^\mathsf T\), is symplectic when
\[
M^\mathsf TJM=J,\qquad
J=\begin{pmatrix}0&1\\-1&0\end{pmatrix}.
\]

For a centered beam with covariance matrix \(\Sigma\), define the
geometric RMS emittance
\[
\varepsilon:=\sqrt{\det\Sigma},
\qquad
\Sigma=
\varepsilon
\begin{pmatrix}
\beta_T&-\alpha_T\\
-\alpha_T&\gamma_T
\end{pmatrix},
\qquad
\beta_T\gamma_T-\alpha_T^2=1.
\]
The dimensions are
\[
[\beta_T]=L,\quad [\alpha_T]=1,\quad
[\gamma_T]=L^{-1},\quad [\varepsilon]=L.
\]

\begin{proposition}[Qualified geometric-emittance invariance]
If \(\Sigma_2=M\Sigma_1M^\mathsf T\) and \(M\) is a \(2\times2\)
symplectic map in the same canonical coordinate convention, then
\(\det\Sigma_2=\det\Sigma_1\) and
\(\varepsilon_2=\varepsilon_1\).
\end{proposition}

\begin{proof}
A \(2\times2\) symplectic matrix has \(\det M=1\). Hence
\(\det\Sigma_2=(\det M)^2\det\Sigma_1=\det\Sigma_1\).
\end{proof}

During ideal acceleration the geometric emittance in trace-space
coordinates undergoes adiabatic damping. The corresponding normalized
emittance is
\[
\varepsilon_n=\beta_{\rm rel}\gamma_{\rm rel}\varepsilon.
\]
Its invariance still requires ideal uncoupled transport and excludes
scattering, radiation, noise, nonlinear filamentation, collective
effects, and losses. The word ``invariant'' without this hypothesis
list is rejected.

\section{Bunches as hidden distributions}

Let \(u=(x,p_x,y,p_y,z,\delta)\) denote a declared six-dimensional
coordinate chart and let
\[
f_t(u)\ge0,\qquad
\int f_t(u)\,\dd\mu_\Gamma(u)=1
\]
be a normalized hidden beam law with a typed phase-space measure.
An intensity density uses a different normalization,
\(\int f_t\,\dd\mu_\Gamma=N_b\).

Profile monitors observe marginals or convolutions, for example
\[
\rho_x^{\rm ideal}(x)
=\int f_t(u)\,\dd\mu_{\Gamma\setminus x},
\qquad
Y_x\sim\Kcal_\lambda\bigl(\rho_x^{\rm ideal}\bigr).
\]
Equal one-dimensional marginals do not identify the joint law.
Different transverse correlations, tails, halo populations, and
loss risks can therefore agree at the profile level.

\section{Collision and detector reconstruction}

For reconstructed final-state energy four-vectors
\(\Pi_i^\mu=(E_i,\bm\pi_i)\),
\[
\mu_{\rm inv}^2
=
\left(\sum_iE_i\right)^2
-\left\lVert\sum_i\bm\pi_i\right\rVert^2.
\]
The same physical invariant in SI momentum coordinates is
\[
m_{\rm inv}^2c^4
=
\left(\sum_iE_i\right)^2
-c^2\left\lVert\sum_i\bm p_i\right\rVert^2.
\]
Detector data do not provide the exact \(\Pi_i^\mu\). Tracks, clusters,
timing, calibration constants, object identification, and selections
enter an estimator
\[
\widehat\mu_{\rm inv}
=\widehat\theta_\lambda(Y),\qquad
\Sigma_{\widehat\mu}
\ \text{or a likelihood interval}.
\]

For polar angle \(\theta\),
\[
p_T=\norm{\bm p_\perp},\qquad
\eta=-\ln\tan\frac{\theta}{2}.
\]
The logarithm is dimensionless. Under a longitudinal boost, \(p_T\)
is unchanged. True rapidity
\[
y=\frac12\ln\frac{E+\pi_z}{E-\pi_z}
\]
shifts additively, so rapidity differences are boost-invariant.
Pseudorapidity equals rapidity only in the massless limit; it is not
an unrestricted Lorentz scalar.

\section{Transformation classes}

\begin{longtable}{@{}L{37mm}L{43mm}L{67mm}@{}}
\toprule
Class & Object & Correct status \\
\midrule
\endhead
Lorentz &
\(\Pi^2,\mu_{\rm inv}\) &
invariant under proper Lorentz transformations \\
Longitudinal boost &
\(p_T,y,\Delta y\) &
\(p_T\) unchanged, \(y\) shifts, \(\Delta y\) invariant \\
Symplectic transport &
\(\varepsilon=\sqrt{\det\Sigma}\) &
invariant under the stated ideal map and coordinate convention \\
Adiabatic acceleration &
\(\varepsilon_n\) &
ideal normalized-emittance invariant under additional hypotheses \\
Unit chart &
physical \(Q\) &
same quantity; numerical coordinate transforms covariantly \\
Detector reconstruction &
\(\widehat\mu_{\rm inv}\) &
estimator with calibration and uncertainty, not an exact invariant \\
Magnet setting &
\(B\rho\) &
frame- and species-dependent rigidity relation, not Lorentz invariant \\
\bottomrule
\end{longtable}

\section{Non-identifiability}

For a deterministic observation \(y=\Ocal_\lambda(x)\), the exact
fiber is
\[
\Fcal_y=\{x:\Ocal_\lambda(x)=y\}.
\]
For a stochastic channel with likelihood \(k_\lambda(y\mid x)\), the
compatibility set becomes
\[
\Fcal_{y,\lambda}
=\{x:k_\lambda(y\mid x)>0\},
\]
or, more usefully, a likelihood or posterior level set. A ``size of
the fiber'' is undefined until a metric, measure, prior, threshold,
and quotient by representation symmetries are specified.

This distinction blocks three invalid inferences:
\begin{enumerate}[label=(\roman*)]
  \item a beam profile does not determine the full phase-space law;
  \item a reconstructed invariant mass does not determine a unique
  partonic history;
  \item an event display is not the collision event itself.
\end{enumerate}

\section{Reproducible Run 3 numerical audit}

The following values use the exact SI definitions of \(c,h,e\), the
CODATA 2022 proton mass energy, the CERN public \(6.8\,\TeV\)
reference, and the rounded public circumference.

\begin{center}
\small
\begin{tabularx}{\linewidth}{@{}Y L{40mm} L{31mm}@{}}
\toprule
Quantity & Value & Status \\
\midrule
Reference proton total energy & \(6800\,\mathrm{GeV}\) & operational reference \\
Proton mass energy & \(0.93827208943\,\mathrm{GeV}\) & CODATA 2022 \\
Lorentz factor \(\gamma\) & \(7.24736467876\times 10^{3}\) & derived \\
Speed ratio \(\beta\) & \(0.99999999048060\) & derived \\
Stable \(1-\beta\) & \(9.519404390069\times 10^{-9}\) & derived \\
Speed deficit \(c-v\) & \(2.853845641\,\mathrm{m\,s^{-1}}\) & derived \\
Momentum & \(6799.999935268\,\mathrm{GeV}/c\) & derived \\
Magnetic rigidity \(B\rho\) & \(22682.358258\,\mathrm{T\,m}\) & derived \\
Revolution frequency & \(11245.450135\,\mathrm{Hz}\) & rounded circumference \\
Symmetric head-on \(\sqrt{s}\) & \(13600\,\mathrm{GeV}\) & derived \\
\bottomrule
\end{tabularx}

\end{center}

The revolution-frequency value has more displayed digits only as a
regression output; its physical precision is limited by the rounded
circumference and the orbit definition. The approximate public bunch
count and population would imply a beam-energy scale of about
\(4.36\times10^8\,\mathrm J\), but this is not an event-by-event or
fill-specific measurement.

\section{Claim audit}

\begin{center}
\begin{tabularx}{\linewidth}{@{}L{49mm}L{27mm}Y@{}}
\toprule
Statement & Status & Required hypotheses or reason \\
\midrule
Mass-shell equivalence across SI/HEP charts & theorem & explicit \(c\)-map and common signature \\
Run 3 numerical audit & derived check & declared source values and epoch \\
Geometric-emittance invariance & proposition & linear symplectic map, fixed canonical convention \\
Normalized-emittance invariance & ideal model & adiabatic acceleration and excluded nonideal effects \\
Equal-time longitudinal contraction & standard kinematics & inertial frames and declared simultaneity slices \\
Penrose--Terrell appearance & standard optical result & past-light-cone image of an extended source \\
Direct optical image of an LHC proton & rejected & no such detector channel is modelled \\
New LHC or particle-physics prediction & not claimed & reference/infrastructure module only \\
\bottomrule
\end{tabularx}
\end{center}

\section*{References}
\addcontentsline{toc}{section}{References}

\begin{enumerate}[label={[\arabic*]}]
  \item BIPM, \emph{The International System of Units (SI)},
  9th ed., version 4.01.
  \url{https://www.bipm.org/documents/20126/41483022/SI-Brochure-9-EN.pdf}
  \item P.\ J.\ Mohr et al., ``CODATA recommended values of the
  fundamental physical constants: 2022,'' 2025.
  \url{https://physics.nist.gov/cuu/pdf/JPCRD2022CODATA.pdf}
  \item Particle Data Group, ``Physical Constants,'' in
  \emph{Review of Particle Physics 2024}.
  \url{https://pdg.lbl.gov/2024/reviews/rpp2024-rev-phys-constants.pdf}
  \item CERN, ``Large Hadron Collider: facts and figures.''
  \url{https://home.cern/science/accelerators/large-hadron-collider/}
  \item CERN, ``CERN bids farewell to the LHC and enters Long
  Shutdown 3,'' 29 June 2026.
  \href{https://home.cern/cern-bids-farewell-to-the-lhc-and-enters-long-shutdown-3/}
  {CERN news page, 29 June 2026}.
  \item O.\ Br\"uning et al. (eds.), \emph{LHC Design Report},
  Vol.\ I, CERN, 2004.
  \url{https://cds.cern.ch/record/782076/files/CERN-2004-003-V1-ft.pdf}
  \item E.\ D.\ Courant and H.\ S.\ Snyder, ``Theory of the
  alternating-gradient synchrotron,'' \emph{Annals of Physics}
  \textbf{3} (1958), 1--48.
  \url{https://doi.org/10.1016/0003-4916(58)90012-5}
  \item R.\ Penrose, ``The apparent shape of a relativistically
  moving sphere,'' \emph{Math. Proc. Camb. Phil. Soc.}
  \textbf{55} (1959), 137--139.
  \url{https://doi.org/10.1017/S0305004100033776}
  \item J.\ Terrell, ``Invisibility of the Lorentz contraction,''
  \emph{Physical Review} \textbf{116} (1959), 1041--1045.
  \url{https://doi.org/10.1103/PhysRev.116.1041}
\end{enumerate}

\end{document}
